\section{Results}
\label{sec:comparison-results}
The measurements support three findings. Placement substantially changes the
typical burst of a fixed implementation. Native kernels change the cost of
obtaining that reduction. And complete host behavior depends on
synchronization, alignment, and release timing as well as on analyzer cost.

\subsection{Placement and native implementation cost}
Table~\ref{tab:study-baselines} reports the scalar reference shape. Matched
distribution reduces the median per-hop peak from \StudyBatchPeak{} to
\StudyCorePeak\,$\mu$s, a factor of \StudyPlacementRatio, while cost rises
from \StudyBatchCost{} to \StudyCoreCost\,ns per engine sample. Across all
three baseline lengths, all four matched implementation families, and both
sessions, the distributed or hybrid path has the lower session hop-peak
summary in \StudyPlacementSigns{} of \StudyPlacementComparisons{} comparisons.
This is descriptive sign consistency over dependent cells; it is neither a
significance test nor a ranking of unrelated kernels.

\input{.build/paper-study/baselines.tex}

Native kernels lower the price of a similar reduction. The PFFFT hybrid costs
\StudyPffftHybridCost\,ns/sample with an \StudyPffftHybridPeak\,$\mu$s hop
peak, against \StudyPffftBatchCost\,ns/sample and \StudyPffftBatchPeak\,$\mu$s
for its matched batch path; vDSP and FFTW supply further tradeoff points.
Figure~\ref{fig:study-tradeoff} shows every process at this shape. A small
peak does not imply minimal aggregate work, and the variation between sessions
is shown rather than folded into a single score per provider. Redistributing
work changes dispatch, vector-call sizes, and locality, so even matched
arithmetic need not produce identical elapsed cost.

\input{.build/paper-study/tradeoff.tex}

At this shape even the scalar batch path uses only a few percent of its
1,333\,$\mu$s block budget. The reference exists to isolate work placement;
its practical significance must be judged at the complete-module and host
boundaries below. Callback p99 and retained maxima also differ markedly from
the median hop peak: the per-hop statistic keeps sparse frame work visible,
whereas maxima retain interruptions and other excursions. None of these
statistics is a safe execution-time bound.

\subsection{Completion horizon and four-channel analysis}
The native horizon sweep makes publication age an explicit design variable
(Table~\ref{tab:study-horizons}). At the demanding controlled shape
$N=16384$, $H=240$, contracting the four-channel vDSP hybrid's horizon from
$H$ to $H/2$ to $H/4$ cuts endpoint age from 4.98 to 2.48 to 1.23\,ms while
raising its median hop peak from \StudyHorizonFullPeak{} to
\StudyHorizonHalfPeak{} to \StudyHorizonQuarterPeak\,$\mu$s. Because the
native FFT remains indivisible, a shorter horizon packs the surrounding stages
more densely without distributing the transform itself. Behavior at small
sizes and for other kernels need not be monotonic; all 39 granularity cells
are retained in the accompanying data.

\input{.build/paper-study/horizons.tex}

At the same four-channel shape, the SIMD production core has a
\StudySimdExtremePeak\,$\mu$s peak at \StudySimdExtremeCost\,ns/sample,
against \StudyHorizonFullPeak\,$\mu$s at \StudyHorizonFullCost\,ns/sample for
the full-hop vDSP hybrid. Here the native hybrid is better on both recorded
cost and typical peak. Channel-level SIMD therefore does not by itself give
the production core an advantage over batched native channels; which approach
wins depends on shape and kernel. These controlled rows omit module
conditioning and display conversion.

\subsection{Shipped defaults and demanding module settings}
Table~\ref{tab:study-modules} measures actual module calls with paced
releases. With its defaults, Fourier's hop peak is
\StudyModuleDefaultPeak\,$\mu$s, \StudyModuleDefaultBudgetPercent\% of the
block budget. At $N=16384$ with a 5\,ms hop, the complete four-port module
reaches \StudyModuleExtremePeak\,$\mu$s, or
\StudyModuleExtremeBudgetPercent\%. Enabling both smoothing modes raises this
to \StudyModuleBothPeak\,$\mu$s (\StudyModuleBothBudgetPercent\%), with
\StudyModuleBothComputeMisses{} compute-budget exceedances and
\StudyModuleBothMisses{} late releases in 3,840 blocks. The hardest shipped
setting is thus a materially different regime from both the scalar reference
and the defaults.

\input{.build/paper-study/modules.tex}

Using one active port does not make Fourier a scalar analyzer: its four-lane
analysis still runs, and additional voices are summed into each port's lane
rather than becoming separate FFT channels. The table therefore captures
control and input costs and does not imply a fourfold saving when three ports
are unpatched. No paced native complete-module comparison was collected at the
extreme shape, and the continuous native core results cannot fill that gap.
These settings expose a limitation of the production path; they are not
evidence that an alternative would eliminate its missed releases.

\subsection{Actual Rack engine, alignment, and missed releases}
Four aligned Fourier modules on one engine thread produce a large difference
in concentration (Table~\ref{tab:study-host}). Production's typical hop peak
of \StudyHostCorePeak\,$\mu$s uses \StudyHostCoreBudgetPercent\% of the
budget, against \StudyHostBatchPeak\,$\mu$s (\StudyHostBatchBudgetPercent\%)
for vDSP batch and \StudyHostHybridPeak\,$\mu$s
(\StudyHostHybridBudgetPercent\%) for the vDSP hybrid. The price is aggregate
block wall cost: \StudyHostCoreCost{} versus \StudyHostBatchCost{} and
\StudyHostHybridCost\,ns/sample, respectively. In this regime the burst
reduction frees a substantial fraction of the block for other work, rather
than shaving a few microseconds from an otherwise idle callback.

\input{.build/paper-study/host.tex}

Four engine threads change the picture. With four aligned modules, the three
paths converge to typical hop peaks of 100--115\,$\mu$s
(Figure~\ref{fig:study-host}), and the empty four-thread engine has a much
larger callback p99 than the one-thread control. Rack's per-sample barriers
and worker dispatch are a plausible mechanism, but these runs do not separate
that cost from CPU placement or power-state effects, and the synchronization
code alone does not establish it. Table~\ref{tab:study-host-accounting} adds
an important observation: the empty four-thread engine consumes about
\StudyEmptyCpuCores{} CPU-seconds per second of nominal audio, consistent with
workers spinning between releases, and its process CPU stays near that level
when analyzers are present. The one-thread engine, by contrast, can sleep
between releases. Moving to four threads therefore changes the activity and
potential power-state regime as well as the available parallelism. We report
this as the observed behavior of the pinned headless engine, not as an
isolated measurement of barrier overhead.

\input{.build/paper-study/host-accounting.tex}

\input{.build/paper-study/host-figure.tex}

With 16 aligned modules on four threads, production's typical peak is
\StudyHostManyCorePeak\,$\mu$s against \StudyHostManyBatchPeak\,$\mu$s for
vDSP batch, and staggering the native instances lowers the batch peak to
\StudyHostStaggerBatchPeak\,$\mu$s. Staggering is thus a practical
alternative whenever frame phases can be chosen. The predeclared sweep of 0,
16, and 64 background modules with four staggered analyzers did not reach a
repeatable saturation point, and its latency summaries are not monotonic in
the number of background modules. We retain those cells but do not treat the
background-module count as a calibrated CPU load.

Release misses call for a stronger qualification. On one thread with four
modules, production and vDSP batch have \StudyHostCoreMisses{} and
\StudyHostBatchMisses{} late releases despite very different typical peaks,
and the vDSP hybrid has \StudyHostHybridMisses. The corresponding FFTW batch
row has \StudyHostFftwBatchMisses, but that provider-specific outcome does not
demonstrate a general benefit of scheduling. With 16 aligned modules,
production has \StudyHostManyCoreMisses{} late releases while vDSP batch has
none in these runs, and even empty-engine controls miss releases. These counts
are observations under synthetic pacing and include wake delay and backlog.
They support comparisons of headroom and concentration; they do not establish
fewer audio underruns or a stable ordering of rare-event reliability.

\subsection{Consumer age and numerical limits}
Table~\ref{tab:study-consumer} reports the headless consumer. At 60\,Hz,
production and the full-hop vDSP hybrid have median endpoint-age bounds of
38--40\,ms, against 8--10\,ms for native batch. The difference, about one
30\,ms hop, matches the publication contract. At 30\,Hz, or with polling
stalls, some complete Fourier snapshots are superseded before they are read:
correct ownership protects the data a consumer holds, but it does not
guarantee that the consumer sees every frame. Logical age bounds do not
measure screen refresh, perceptual acceptability, or real-time freshness
during a backlog.

\input{.build/paper-study/consumer.tex}

All numerical replays pass their declared policies
(Table~\ref{tab:study-numerical}), including non-power-of-two hops, high
overlap, weak and decaying input, silence, and complete-module lifecycle
changes. Preparing the study exposed premature underflow in naively squared
magnitudes, and the measured source uses scaled magnitudes in both production
and native paths. Long decays still reach the flagged FTZ/display tail, where
observed relative error can reach one, so the ordinary relative budgets do not
extend to that region. Retaining the tail with its absolute diagnostics is
more informative than labeling the entire decay uniformly accurate.

\input{.build/paper-study/numerical.tex}

Lifecycle replays check freeze and resume, window and band changes, geometry
changes, reset, and sample-rate changes separately from the stationary
comparisons. These mutations occur between engine blocks, and their durations
and allocation probes are recorded outside ordinary block intervals. The study
thus verifies serialized lifecycle behavior; it does not show that arbitrary
concurrent control changes, or all host callbacks, are allocation-free.
